The usefulness of the concept of weak-* convergence is to a large extent based on the following result.

\begin{Theorem}{Weak compactness, Alaoglu}
\label{thm:alaoglu-weak-star-subsequence}
Let $X$ be a separable Banach space and let $f_n$ be a bounded sequence in $X^*$. Then $f_n$ has a weakly-* convergent subsequence.
\end{Theorem}

\begin{Proof}
Let $x_k$, $k\in\mathbb{N}$ be a sequence that is dense in $X$. We first note that if $f_n(x_k)$ converges for every $k\in\mathbb{N}$, then $f_n(x)$ converges for every $x\in X$. To see this, let $M$ be an upper bound for $\|f_n\|$. For given $\varepsilon>0$ and $x\in X$, we can choose $k$ such that $\|x-x_k\|\leq\varepsilon/(3M)$. We then have
\begin{equation}\label{eq:alaoglu-dense-set-cauchy-estimate}
\begin{aligned}
|f_n(x)-f_m(x)|&\leq|f_n(x)-f_n(x_k)|+|f_n(x_k)-f_m(x_k)|\\
&\qquad+|f_m(x_k)-f_m(x)|\\
&\leq\dfrac{2}{3}\varepsilon+|f_n(x_k)-f_m(x_k)|,
\end{aligned}
\end{equation}
and the last term on the right-hand side is less than $\varepsilon/3$ if $m$ and $n$ are large enough. Hence $f_n(x)$ converges for every $x\in X$. The rest of the proof is a standard diagonal argument. Since $f_n(x_1)$ is bounded, we can extract a subsequence $f_{n1}$ such that $f_{n1}(x_1)$ converges. From this subsequence, we extract another subsequence $f_{n2}$ such that $f_{n2}(x_2)$ converges. We proceed like this inductively. The diagonal sequence $f_{nn}$ has the property that $f_{nn}(x_k)$ converges for every $k$.
\end{Proof}

\begin{Remark}
\label{rem:weak-compactness-approximate-pde-solutions}
In the typical applications of this theorem, the $f_n$ are approximate solutions to a PDE. While it may be difficult to show that the sequence converges, one can often obtain bounds on the sequence which imply a weakly convergent subsequence. The limit of the subsequence is then shown to be the solution of the original PDE.
\end{Remark}
